\section{Physical realization and transfer at an exterior maximum}\label{sec:transfer}

Two different precisions of matching enter the proof. At the inner scale it
suffices to decide which of two exponentially small gaps is larger. At the
exterior maximum two input ordinates a bounded physical distance apart
have the same limiting scaled height. To distinguish their upper hits we
first match the references and then compare ratios of gaps on each side.
The estimates in~\eqref{eq:log-gap-quad} and~\eqref{eq:cap-fine} give these
two comparisons.

\subsection{Physical realization with one common reference over a quadratic germ}\label{subsec:quad-realize}
Let $G_s(z)=E(z)+sC(z)$ ($s=\pm 1$, $z\ge 0$) be fixed polynomial profiles with $E$ even ($\deg E=m\ge 8$, leading coefficient $c>0$), $C$ odd ($\deg C=j<m$), $G_s(z)>0$ for all $z>0$, and low-order germs
\begin{equation}\label{eq:quad-germs}
  E(z) = \eta z^2 + O(z^4),\qquad C(z) = \gamma_3 \eta z^3 + O(z^5)\qquad (\eta>0,\ \gamma_3>0\text{ fixed}).
\end{equation}
The quadratic coefficient $\eta$ stays fixed as the spatial scale $L$ increases.
Assume that on each side $s\in\{\pm 1\}$ independently:
\begin{enumerate}
  \item[\textup{(H1)}] \emph{Nondegenerate ordered fold windows.} The positive critical points of $G_s$ on $(0,\infty)$ consist of $K_s\ge 0$ nondegenerate local maxima $p_{s,l}$ and subsequent positive local minima $q_{s,l}$ ($0=:q_{s,0}<p_{s,1}<q_{s,1}<\dots<p_{s,K_s}<q_{s,K_s}$) satisfying $G_s''(p_{s,l})<0<G_s''(q_{s,l})$ and $G_s'(z)>0$ on $(0,\infty)\setminus\bigcup_{l=1}^{K_s}[p_{s,l},q_{s,l}]$, whose critical height windows $[\mu_{s,l},M_{s,l}]=[G_s(q_{s,l}),G_s(p_{s,l})]\subset(0,\infty)$ are pairwise disjoint and strictly increasing: $0=:\mu_{s,0}<\mu_{s,1}<M_{s,1}<\dots<\mu_{s,K_s}<M_{s,K_s}$.
  \item[\textup{(H2)}] \emph{Transverse inner landings.} Because $G_s$ increases strictly on $[q_{s,l-1},p_{s,l}]$ from $\mu_{s,l-1}<\mu_{s,l}$ to $M_{s,l}>\mu_{s,l}$, each minimum level $\mu_{s,l}$ ($1\le l\le K_s$) admits a unique inner landing $r_{s,l}\in(q_{s,l-1},p_{s,l})$ solving $G_s(r_{s,l})=\mu_{s,l}$, which is transverse ($G_s'(r_{s,l})>0$) with $G_s(z)>\mu_{s,l}$ throughout $(r_{s,l},q_{s,l})\cup(q_{s,l},\infty)$.
\end{enumerate}

Write $\mathcal{C}_s=\{\mu_{s,l},M_{s,l}:1\le l\le K_s\}$ for the finite set of critical fold-endpoint heights on side $s$, and put $\mathcal{C}=\mathcal{C}_+\cup\mathcal{C}_-\subset(0,\infty)$. A positive height $Y>0$ is \emph{regular} when $Y\in(0,\infty)\setminus\mathcal{C}$—that is, away from the fold endpoints $\mu_{s,l},M_{s,l}$, \emph{including} interior window levels $\mu_{s,l}<Y<M_{s,l}$. For \emph{every} regular height $Y\in(0,\infty)\setminus\mathcal{C}_s$:
\begin{itemize}
  \item The largest positive root $v_s(Y)$ of $G_s(z)=Y$ lies outside $\bigcup_{l=1}^{K_s}[p_{s,l},q_{s,l}]$, so $G_s'(v_s(Y))>0$ and $G_s(z)>Y$ on $(v_s(Y),\infty)$ (transverse inward capture from $z=+\infty$).
  \item The first positive root $r_s^{\mathrm{in}}(Y)=\min\{z>0:G_s(z)=Y\}$ satisfies $G_s(z)<Y$ on $[0,r_s^{\mathrm{in}}(Y))$ and $G_s'(r_s^{\mathrm{in}}(Y))>0$ (transverse outward capture from $z=0$): if $\mu_{s,l_0}<Y<M_{s,l_0}$ for some (unique) $l_0$, then $r_s^{\mathrm{in}}(Y)\in(r_{s,l_0},p_{s,l_0})$ lies on the inner rising branch above the transverse landing $r_{s,l_0}$, whereas if $Y\notin\bigcup_{l=1}^{K_s}[\mu_{s,l},M_{s,l}]$ then $r_s^{\mathrm{in}}(Y)=v_s(Y)$.
\end{itemize}
Accordingly, for every $Y\in(0,\infty)\setminus\mathcal{C}$ we define \emph{both} the incoming and upper retained actions by
\begin{align}
  J_s^I(Y) &= \int_0^{v_s(Y)}\frac{G_s'(z)^2}{z}\,\dd z - \sum_{l:\,\mu_{s,l}<Y}\int_{r_{s,l}}^{q_{s,l}}\frac{G_s'(z)^2}{z}\,\dd z,\label{eq:JI-def}\\
  J_s^U(Y) &= \begin{cases}
    \displaystyle\int_0^{r_{s,l_0}}\frac{G_s'(z)^2}{z}\,\dd z - \sum_{l<l_0}\int_{r_{s,l}}^{q_{s,l}}\frac{G_s'(z)^2}{z}\,\dd z, & \text{if }\mu_{s,l_0}<Y<M_{s,l_0},\\[1.1ex]
    J_s^I(Y), & \text{if }Y\notin\bigcup_{l=1}^{K_s}[\mu_{s,l},M_{s,l}].
  \end{cases}\label{eq:JU-def}
\end{align}
Inside a fold window the launch and the reference remain separated
until the reference jumps to the inner landing $r_{s,l_0}$. Their gap
contracts exponentially along the common part of the two paths below
the jump. This is why the upper limit in~\eqref{eq:JU-def} is the same
for all launch heights within that window.
For negative entrance magnitude $w>0$, define the absolutely convergent integral
\begin{equation}\label{eq:N-def}
  N(w) = \int_0^\infty z\left(\frac{1}{w+G_+(z)} - \frac{1}{w+G_-(z)}\right)\dd z.
\end{equation}

Fix three separate finite sample sets: an upper-launch set $\mathcal{Y}_U\subset(0,\infty)\setminus\mathcal{C}$ on which $J_+^U(Y)-J_-^U(Y)\neq 0$ (which may include levels inside open fold windows $(\mu_{s,l},M_{s,l})$), a positive-incoming set $\mathcal{Y}_I\subset(0,\infty)\setminus\mathcal{C}$ on which $J_-^I(Y)-J_+^I(Y)\neq 0$, and a negative-incoming set $\mathcal{W}_N\subset(0,\infty)$ on which $N(w)\neq 0$. Choose a reference level $T>\max(\mathcal{Y}_U\cup\mathcal{Y}_I\cup\mathcal{C})$ ($\max\varnothing:=0$), and define the physical polynomial family at scale $L\gg 1$:
\begin{equation}\label{eq:quad-field}
  S = \frac{L^2}{\eta},\qquad \eps = \frac{L^2}{S^2} = \frac{\eta^2}{L^2},\qquad
  F_{L,a}(x) = S\big(E(x/L)+C(x/L)\big) + ax,\quad a\in[-1,0].
\end{equation}

\begin{proposition}[Simultaneous physical realization over a quadratic core]\label{prop:physical-quad}
For all sufficiently large $L$:
\begin{enumerate}
  \item[\textup{(i)}] Every upper launch $h>0$ and incoming ordinate $y\in\R$ has transverse lower-axis hits $w_s(h)$ and $I_s(y)$ on both sides $s=\pm 1$ for all $a\in[-1,0]$, there exists an interior parameter $a_L\in(-1,0)$ closing the common upper reference $w_+(ST)=w_-(ST)=:w_*$, and whenever $a\in(-1,0)$ one has $\D(h_*)>0$ at sufficiently small positive $h_*>0$.
  \item[\textup{(ii)}] Uniformly in $a\in[-1,0]$, the positive endpoint gaps $g_s^U=w_s(ST)-w_s(SY)>0$ ($Y\in\mathcal{Y}_U$) and $g_s^I=I_s(SY)-w_s(ST)>0$ ($Y\in\mathcal{Y}_I$) satisfy
  \begin{equation}\label{eq:log-gap-quad}
    -\eps\log g_s^U \to J_s^U(Y)\ (Y\in\mathcal{Y}_U),\quad
    -\eps\log g_s^I \to J_s^I(Y)\ (Y\in\mathcal{Y}_I)\quad (L\to\infty).
  \end{equation}
  At $a=a_L$, $\D(SY)=g_-^U-g_+^U$ has the strict sign of $J_+^U(Y)-J_-^U(Y)$ for every $Y\in\mathcal{Y}_U$, and $I_+(SY)-I_-(SY)=g_+^I-g_-^I$ has the strict sign of $J_-^I(Y)-J_+^I(Y)$ for every $Y\in\mathcal{Y}_I$.
  \item[\textup{(iii)}] Uniformly in $a\in[-1,0]$ at each $w\in\mathcal{W}_N$, $(I_+(-Sw)-I_-(-Sw))/(S\eps)\to N(w)$ as $L\to\infty$, hence $I_+(-Sw)-I_-(-Sw)$ has the strict sign of $N(w)$.
\end{enumerate}
\end{proposition}

\begin{proof}
(i)~Since $G_s(z)>0$ for $z>0$, $G_s(z)\sim\eta z^2$ near $0$, and $G_s(z)\sim c z^m$ at $\infty$, there exists $c_0>0$ such that $G_s(z)\ge c_0 z^2$ for all $z\ge 0$. In physical coordinates $F_{L,a}(sx)=SG_s(x/L)+sax\ge (c_0/\eta)x^2-|x|\to+\infty$ as $x\to+\infty$ and $|F_{L,a}'(0)|=|a|\le 1<2$, so Lemmas~\ref{lem:half-returns}(iii) and~\ref{lem:incoming-rays}(i) guarantee global transverse half-returns $w_s(h)$ and $I_s(y)$ for all $a\in[-1,0]$. Moreover, when $a\in(-1,0)$, $F_{L,a}(x)-F_{L,a}(-x)=2ax+O_L(x^3)<0$ on $(0,x_0]$ for small $x_0>0$, so along~\eqref{eq:scalar-u} any small launch $h_*>0$ remaining in $x\le x_0$ satisfies $v_+(y)>v_-(y)$ down to $\Sigma_-$ and hence $\D(h_*)>0$.

Write $v=x^2/2$ in physical units, so $\dd v/\dd(-y)=y-F_{L,a}(s\sqrt{2v})\le y-c_1 v+c_2$ with $c_1,c_2>0$ independent of $L,a$. Integrating downwards from $(0,ST)$ gives $v(y)\le C_1(1+y)$ for all $0\le y\le ST$ (and similarly for any arc entering at height $O(S)$ with $v=O(L^2)$, since $L^2 e^{-c_1\Delta y}$ decays exponentially over physical descent). On bounded physical sets,~\eqref{eq:quad-germs} and $S=L^2/\eta$ give $F_{L,a}(sx)\to x^2+sax$ in $C^2$. Hence for $y\ge 1$ and sufficiently small $v\le v_{\min}$, $\dd v/\dd(-y)\ge 1/2$, which drives every arriving arc into a fixed compact interval $[v_0,V_0]\subset(0,\infty)$ throughout $1\le y\le 8$, and the section $y=8$ is transverse since $\dot y=-x<0$ there.
For the limiting fields $x^2+sax$ the incoming and terminal sections
are both transverse. The derivative between them is continuous and
nonzero on the compact set of section data and parameters. Compact
flow dependence gives the same bounds for the physical fields. Thus the terminal-hit map $v(8)\mapsto B_s(\sqrt{2v(8)},8)$ has derivative bounded above and away from zero uniformly for large $L$ and $a\in[-1,0]$.

To bracket the reference match at $ST$, set $Y_L=4\log L$. On $1\le y\le Y_L$, $v_\pm(y)\le C_1(1+Y_L)$ gives $x_\pm=O(\sqrt{\log L})$, where~\eqref{eq:quad-germs} yields $\partial_v(SE(x/L))=2+O((\log L)/L^2)$ and $SC(x/L)=\gamma_3 x^3/L+O((\log L)^{5/2}/L^3)$. Subtracting~\eqref{eq:scalar-u} for $s=+1$ and $s=-1$, the difference $d(y)=v_+(y)-v_-(y)$ satisfies
\[
  \frac{\dd d}{\dd(-y)} = -k(y)d - \big(SC(x_+/L)+ax_+\big) - \big(SC(x_-/L)+ax_-\big),\qquad 1\le k(y)\le 3
\]
for large $L$. Integrating from $y=Y_L$ down to $y=1$, the homogeneous memory $d(Y_L)\exp(-\int_1^{Y_L}k)$ is $O((\log L)e^{-Y_L})=O((\log L)L^{-4})$. At $a=0$, the forcing $-SC(x_+/L)-SC(x_-/L)$ is nonpositive on $[1,Y_L]$ and at most $-c_3/L$ throughout $y\in[1,2]$ (where $x_\pm\ge\sqrt{2v_0}>0$), dominating $O((\log L)L^{-4})$ and forcing $d(1)<0$. Since $SC(x/L)$ has positive cubic germ near zero and coercivity keeps $v_\pm(y)\le 2C_1$ bounded on $y\le 1$ all the way to the lower-axis hit, $x_\pm/L=O(1/L)$ stays strictly inside the positive-germ neighborhood of $C$ down to the hit. Likewise at $a=-1$, $ax+\gamma_3 x^3/L\le -\tfrac12 x<0$ on the entire bounded range of $x_\pm$ below $Y_L$, giving forcing at least $+c_4>0$ on $[1,2]$ and $d(1)>0$, with the opposite vector-field inequality holding down to the hit. Trajectory comparison on a single side therefore gives $w_+(ST)<w_-(ST)$ at $a=0$ and $w_+(ST)>w_-(ST)$ at $a=-1$, so continuity in $a$ yields $a_L\in(-1,0)$ with $w_+(ST)=w_-(ST)$.

(ii)~In normalized coordinates $z=x/L$, $u=z^2/2$, $t=y/S$, equation~\eqref{eq:scalar-u} reads
\[
  \eps\,\frac{\dd u}{\dd(-t)} = t - P_{s,L}(u),\qquad P_{s,L}(u)=G_s(\sqrt{2u}) + \frac{sa\eta\sqrt{2u}}{L}.
\]
Fix a section $z_b>v_s(T)$ ($u_b=z_b^2/2$), so the reference $u_T(t)$ from $(0,T)$ stays strictly left of $u_b$ for all $t\le T$. For $Y\in\mathcal{Y}_I$, integrating $\dd t/\dd z=\eps z/(P_{s,L}-t)>0$ inward from $+\infty$ to $z_b$ brings the incoming ray to $(u_b,t_b)$ at $t_b=Y-O(\eps)$, where $u_{\mathrm{in}}(t_b)-u_T(t_b)=u_b-v_s(Y)^2/2+O(L^{-1}+\eps)>0$ is bounded above and away from zero; for $Y\in\mathcal{Y}_U$, the upper launch starts at $(0,Y)$ with $u_T(Y)-0=v_s(Y)^2/2+O(L^{-1}+\eps)>0$ bounded above and away from zero. Hence $g_s^U,g_s^I>0$, and any two ordered solutions $u_1>u_2$ satisfy
\begin{equation}\label{eq:secant-exact}
  \eps\,\frac{\dd}{\dd t}\log(u_1-u_2) = [P_{s,L}]_{u_1,u_2} := \frac{P_{s,L}(u_1)-P_{s,L}(u_2)}{u_1-u_2}.
\end{equation}
Choose $\rho>0$ smaller than half the minimum separation of the distinct
heights in $\{0,Y,T\}\cup\mathcal C$, and remove the scaled height bands $\bigcup_{c_*\in\mathcal{C}_s}[c_*-\rho,c_*+\rho]\cup[0,\rho]$.
On the compact complement of these bands, each tracked rising branch $u_s^*(t)$ of $G_s(\sqrt{2u})=t$ satisfies $\partial_u G_s\ge c_\rho>0$ and $|\dd u_s^*/\dd t|\le C_\rho$. Any narrow fixed tube $|u-u_s^*(t)|\le\nu$ has inward restoring speed $\ge\tfrac12 c_\rho\nu/\eps\gg C_\rho$ as $L\to\infty$: for $Y\in\mathcal{Y}_I$, $G_s>Y$ on $(v_s(Y),z_b]$ and $G_s'(v_s(Y))>0$ bring $u_{\mathrm{in}}$ from $u_b$ into the $v_s(t)^2/2$-tube within $O(\eps)$ descent (including when $\mu_{s,l_0}<Y<M_{s,l_0}$, where $v_s(t)>q_{s,l_0}$ tracks the outer branch smoothly down to $\mu_{s,l_0}$ and jumps with $u_T$ to the transverse landing $r_{s,l_0}$); for $Y\in\mathcal{Y}_U$ outside $\bigcup_l[\mu_{s,l},M_{s,l}]$, $G_s<Y$ on $[0,v_s(Y))$ brings $u_Y$ into the same $v_s(t)^2/2$-tube within $O(\eps)$ descent (with $u_T-u_Y$ bounded away from zero near $u_Y=0$). If $Y\in\mathcal{Y}_U$ satisfies $\mu_{s,l_0}<Y<M_{s,l_0}$, then by~(H1)--(H2) $G_s<Y$ on $[0,r_s^{\mathrm{in}}(Y))$ and $G_s'>0$ throughout $[r_{s,l_0},p_{s,l_0})$, so $u_Y$ enters the inner branch $(r_{s,l_0},p_{s,l_0})$ in $O(\eps)$ descent and tracks it past $\mu_{s,l_0}$ while $u_T$ tracks the outer branch down to $\mu_{s,l_0}$; at $t=\mu_{s,l_0}+\rho$, $u_T-u_Y\to(q_{s,l_0}^2-r_{s,l_0}^2)/2+O(\sqrt\rho)>0$ is bounded above and away from zero, so starting~\eqref{eq:secant-exact} at $t=\mu_{s,l_0}+\rho$ and entering the common landing tube around $r_{s,l_0}$ by $\mu_{s,l_0}-\rho$ omits the $[\mu_{s,l_0},M_{s,l_0}]$ outer segment and skips above $Y$, matching~\eqref{eq:JU-def}. The transition through each minimum takes place inside its excised
band. At $t=\mu_{s,l}-\rho/2$,
the vector field $\eps\,\dd u/\dd(-t)=t-P_{s,L}(u)$ points left with
magnitude bounded below, for fixed $\rho$, between small neighborhoods
of the landing branch and the old minimum, including the minimum itself.
Indeed the limiting profile on this interval is at least $\mu_{s,l}$
between the landing and the minimum, and the landing is transverse.
The same bound holds a little beyond the minimum. An arc still on the
right traverses this interval in $O_\rho(\eps)$ height and enters a
fixed tube about the landing branch before $\mu_{s,l}-\rho$. An arc
already in the tube stays there. This argument applies at each of the
finitely many ordered minima.

On $[8/S,\rho]$, $v\ge v_0$ gives $z\ge\sqrt{2v_0}/L$, keeping $(sa\eta/L)/\sqrt{2u}=O(1)$, so $[8/S,\rho]$ and the $O(\rho)$-width excised bands contribute $O(\rho)$ as $L\to\infty$, while conversion from $v_1-v_2=L^2(u_1-u_2)$ at $y=8$ to $g_s^X$ costs $O(\eps\log L)=o(1)$. On the common tubes $[P_{s,L}]_{u_1,u_2}\to\partial_u G_s(u_s^*(t))$; changing variables $t=G_s(z)$ ($\partial_u G_s\,\dd t=G_s'(z)^2\,\dd z/z$) and letting $L\to\infty$ then $\rho\downarrow 0$ proves~\eqref{eq:log-gap-quad}. At $w_+(ST)=w_-(ST)$, $\D(SY)=g_-^U-g_+^U$ and $I_+(SY)-I_-(SY)=g_+^I-g_-^I$, giving the stated signs.

(iii)~At limiting ordinate $-Sw$ ($w\in\mathcal{W}_N$ fixed), write the inward scaled trajectory as $t_s(z)=-w-\eps V_s(z)$. Then~\eqref{eq:spatial-ode} becomes the integral fixed-point equation
\[
  V_s(z) = \int_z^\infty \frac{r\,\dd r}{w + G_s(r) + sa\eta r/L + \eps V_s(r)}.
\]
Since $G_s(r)\ge 0$ and $G_s(r)\sim c r^m$ ($m\ge 8$), $w+G_s(r)+sa\eta r/L\ge c_w(1+r^m)$ on $[0,\infty)$ for large $L$. On the complete metric space $\{V\in C_b([0,\infty)):0\le V\le M_w\}$ ($M_w=\int_0^\infty r/(c_w(1+r^m))\,\dd r$), the integral operator is a uniform contraction of Lipschitz constant $O(\eps)$ for large $L$, with $V_s(0)=\int_0^\infty r/(w+G_s(r))\,\dd r+O(1/L+\eps)$ uniformly in $a\in[-1,0]$. Since $I_s(-Sw)=-S t_s(0)=S(w+\eps V_s(0))$, subtracting $s=\pm 1$ and dividing by $S\eps$ gives $V_+(0)-V_-(0)\to N(w)$.
\end{proof}

\begin{lemma}[Unbounded-branch tail action and perturbation persistence]\label{lem:tail-action}
Let $G_s=E+sC$ satisfy the hypotheses of Subsection~\ref{subsec:quad-realize} with leading terms $E(z)=cz^m+\dots$ ($c>0$) and $C(z)=bz^j+\dots$ ($j<m$, $b\neq 0$, $j\neq m-2$).
\begin{enumerate}
  \item[\textup{(i)}] As $R\to+\infty$ at regular heights $Y=E(R)$ above all critical levels,
  \begin{equation}\label{eq:tail-action}
    J_+^U(Y)-J_-^U(Y) = J_+^I(Y)-J_-^I(Y) = 2mcb\,\frac{j-m+2}{m+j-2}\,R^{m+j-2} + o\!\left(R^{m+j-2}\right),
  \end{equation}
  and as $w\to+\infty$, $w^{2-(j+2)/m}N(w)\to -2b\int_0^\infty v^{j+1}/(1+cv^m)^2\,\dd v$, so $\sgn N(w)=-\sgn b$.
  \item[\textup{(ii)}] On already fixed finite sample sets $\mathcal{Y}_U,\mathcal{Y}_I\subset(0,\infty)\setminus\mathcal{C}$ and $\mathcal{W}_N\subset(0,\infty)$, every strict sign of $J_+^U-J_-^U$, $J_-^I-J_+^I$, and $N$ persists (with $\mathcal{Y}_U,\mathcal{Y}_I$ remaining disjoint from the perturbed critical set) when $C$ is perturbed by a sufficiently small odd term of degree in $[5,m-1]$, or when $E$ is perturbed by $\mu z^{m+2}$ for sufficiently small $\mu>0$.
\end{enumerate}
\end{lemma}

\begin{proof}
(i)~At $Y=E(R)$ ($R\to\infty$), $G_s(v_s)=E(R)$ gives $v_s(Y)=R-s\frac{b}{mc}R^{j-m+1}(1+o(1))$. On $[0,R]$, $(G_+')^2-(G_-')^2=4E'C'$ gives $\int_0^R 4E'(z)C'(z)\,\dd z/z = \frac{4mcjb}{m+j-2}R^{m+j-2}+o(R^{m+j-2})$. Adding the endpoint displacements $\int_R^{v_+} - \int_R^{v_-} = -2\frac{b}{mc}R^{j-m+1}\cdot\frac{(mcR^{m-1})^2}{R}(1+o(1))=-2mcb R^{m+j-2}+o(R^{m+j-2})$ and noting that the finitely many jump deletions are $O(1)$ gives $\big(\frac{4j}{m+j-2}-2\big)mcb=\frac{2(j-m+2)}{m+j-2}mcb$, proving~\eqref{eq:tail-action}. Rescaling $z=w^{1/m}v$ in~\eqref{eq:N-def} with $w+G_s(z)\ge c_1(w+z^m)$ gives the $N(w)$ limit by dominated convergence since $j<m\implies j+1<2m-1$.

(ii)~Adding tiny higher-degree terms of order $\ge 5$ preserves the low-order germs~\eqref{eq:quad-germs}, simple folds, and transverse landings on compact sets, while a positive $\mu z^{m+2}$ or a sufficiently small odd term of degree $<m$ preserves $G_s'(z)>0$ on the tail and the uniform lower bound $w+G_s(z)\ge c_1(w+z^m)$, yielding continuity of $N(w)$ at each fixed $w$ by dominated convergence.
\end{proof}

\subsection{Matched-maximum transfer at two or three degrees of cost}\label{subsec:cap-transfer}
The input at this stage is a polynomial and a list of actual incoming
half-orbits. Proposition~\ref{prop:physical-quad} is one way to obtain
them. Its separated folds and transverse landings make the singular
action signs into physical comparisons. The transfer itself takes the
transverse hits and their strict signs as data. An exterior maximum
then has two jobs. Its reference heights on the two sides must agree,
and its action difference must match one translated entrance pair.
With those coefficients fixed, adjusting the opposite physical entrance
ordinate by $o(1)$ gives each listed pair a common upper crossing and
preserves its strict comparison.
When $j=m-3$, the even displacement of order $\tau^2$ precedes the odd
perturbation of order $\tau^3$, so the comparison uses the perturbed
\emph{common even} profile. At the finer physical scale we instead use
the common terminal map separately on each side.

\begin{proposition}[Matched-maximum exterior transfer]\label{prop:cap-transfer}
Let $P\in\R[x]$ have $P(0)=0$, even degree $m=\deg P\ge 6$ with leading coefficient $c=[x^m]P>0$, and highest odd power $j=m-g$ with $g\in\{1,3\}$ and $b=[x^j]P\neq 0$. Suppose $P$ possesses $q\ge 0$ hyperbolic limit cycles, and at a nonempty strictly decreasing finite list of ordinates $y_0>y_1>\dots>y_r$ and a real shift $\Delta\in\R$, both incoming trajectories in $I_{+,P}(y_i)-I_{-,P}(y_i-\Delta)$ exist, hit $\Sigma_-$ transversely, and exhibit $r$ strict sign changes ending with $\sigma=\sgn\big(I_{+,P}(y_r)-I_{-,P}(y_r-\Delta)\big)\in\{\pm 1\}$. Put $p=2-g\in\{+1,-1\}$. Then there is a nonempty open subset of exact-degree-$\max(m+2,j+4)$ polynomials (including rational polynomials) on which every field~\eqref{eq:lienard} has at least
\begin{equation}\label{eq:cap-count}
  q + r + 1 + \ind_{\{\sigma\neq\sgn(bp)\}}
\end{equation}
simultaneous hyperbolic limit cycles, retaining all $q$ compact cycles while placing the new upper crossings near $R^m\frac{2c}{m+2}\to+\infty$.
\end{proposition}

The following estimate measures the effect of the part of a passage shared
by all launches. Write $x_+=\max(x,0)$ and $\operatorname{osc}_B f
=\sup_B f-\inf_B f$.

\begin{lemma}[Distortion of a common terminal section]\label{lem:terminal-distortion}
Consider $\epsilon U'=f(t,U)$, with $t$ increasing and $0<\epsilon\le1$.
A compact interval $B\subset(u_c,\infty)$ on $t=b$ flows to its first
downward crossing of $U=u_c>0$
and then to $U=0$. Denote these crossing times by $T(\xi)$ and $K(\xi)$
for $U(b)=\xi\in B$. Suppose the following bounds are uniform, including
over any family of parameters under consideration.
\begin{enumerate}
\item Before the crossing $f$ is piecewise continuous in $t$ and $C^2$
in $U$. The elapsed time and $|f|,|f_U|,|f_{UU}|$ are bounded on
the swept intervals.
In a fixed collar of $u_c$ the field satisfies $f\le-\nu<0$ and
$|f_t|$ is bounded. The smooth flow can be taken piecewise in $t$
at fixed times outside the crossing collar.
\item On the exit strip $0\le U\le u_c$, $u_c$ is bounded,
$f\le-\nu$ and $f,f_t,f_{tt}$ are bounded and continuous in the
inverse chart. Bounds in the collar and strip hold between the arcs.
\end{enumerate}
Let $a\ge0$ be an integrable majorant of $(f_U)_+$ between any two
arcs still above $u_c$, and set $D=\int_b^{\max T}a(t)\,\dd t$.
If $w=|B|e^{D/\epsilon}$ is smaller than the uniform width of the
collar above $u_c$, then $K'>0$ and
\begin{equation}\label{eq:terminal-distortion}
 \operatorname{osc}_{B}(\epsilon\log K')\le Cw .
\end{equation}
In particular if $|B|\le C_0e^{-A/\epsilon}$ and $A-D\ge\gamma>0$,
the right side is $O(e^{-\gamma/\epsilon})$.
For $d_i=|\xi_*-\xi_i|>0$ and $g_i=|K(\xi_*)-K(\xi_i)|$ with
$\xi_*,\xi_1,\xi_2\in B$, this gives
\[
 \left|\epsilon\log(g_1/g_2)-\epsilon\log(d_1/d_2)\right|\le Cw .
\]
\end{lemma}

\begin{proof}
At equal times before the earlier crossing, scalar variation of constants
bounds the distance of two arcs by $w$. At that crossing the other arc
is at most $w$ above $u_c$. It crosses within $\epsilon w/\nu$, by
the collar inequality. Implicit differentiation gives
\[
 T'(\xi)=\frac{-\epsilon}{f(T(\xi),u_c)}
 \exp\!\left(\epsilon^{-1}\int_b^{T(\xi)}f_U(t,U(t;\xi))\,\dd t\right).
\]
On the common time interval the two integrals differ by at most
$\|f_{UU}\|_\infty (\max T-b)w$. Their unmatched ends cost
$\|f_U\|_\infty\epsilon w/\nu$. The log of the section factor has
oscillation at most $\|f_t\|_\infty\epsilon w/\nu^2$.

On the last strip put $h(U,t)=1/f(t,U)$ and solve
$t_U=\epsilon h(U,t)$ from $u_c$ to zero. The inverse flow $L$ has
$\log L'(T)=\epsilon\int_{u_c}^0 h_t(U,t(U;T))\,\dd U$.
Both $h_t$ and $h_{tt}$ are bounded. Two heights in this chart differ
by at most $e^{\epsilon u_c\|h_t\|_\infty}|T_1-T_2|$, so this last
logarithm is Lipschitz in $T$ with constant $O(\epsilon)$.
Composing $K=L\circ T$ proves~\eqref{eq:terminal-distortion}.
Each secant slope of $K$ lies between the extreme values of $K'$ on $B$,
which proves the gap estimate. The inverse-chart proof uses only
continuity in $U$ at the final section.
\end{proof}

In particular, to read a physical scale $\delta$ in the scaled
logarithm of a ratio it suffices that both the entry-action error
and $|B|e^{D/\epsilon}$ are $o(\delta)$. In the exterior transfer
$\delta=\tau^m$ and the entry error is $O(\epsilon|\log\epsilon|)$.
The estimate compares secants from one reference even when the absolute
derivative of the terminal map is itself exponentially small.


\begin{proof}[Proof of Proposition~\ref{prop:cap-transfer}]
Let $R\gg 1$ be the outer spatial scale, $\tau=1/R$, and $\epsilon=\tau^{2m-2}$. Define
\begin{align}
  E_c(z) &= c\!\left(z^m - \frac{m}{m+2}z^{m+2}\right),\qquad A_c(z)=\frac{E_c'(z)}{z}=mc z^{m-2}(1-z^2),\nonumber\\
  G(z)   &= b z^j(1-z^2) + \lambda z^{j+2}(1-z^2) + \theta z^{j+2},\nonumber\\
  F_R(x) &= R^m E_c(x/R) + R^j G(x/R) + P(x) - cx^m - bx^j.\label{eq:cap-poly}
\end{align}
Because $R^m(x/R)^m=x^m$ and $R^j(x/R)^j=x^j$ cancel $-cx^m-bx^j$, $F_R\to P$ in $C^2$ on compact subsets of $\R$ uniformly for bounded $(\lambda,\theta)$, so all $q$ hyperbolic limit cycles of $P$ persist inside a fixed compact disk $K_0\subset\R^2$ for large $R$. On $z>0$, $E_c$ has a unique maximum $E_c(1)=\frac{2c}{m+2}>0$ at $z=1$ and a strictly falling root $v(Y)>1$ for every $Y<E_c(1)$, with $v(0)=\sqrt{(m+2)/m}$. Set $J_o(Y)=\int_1^{v(Y)}E_c'(z)^2\,\dd z/z$ (so $J_o'(Y)=A_c(v(Y))<0$) and $V(Y)=\int_1^{v(Y)}(A_c G'-A_c' G)\,\dd z$. At $\theta=0$, $A_c G'-A_c' G=A_c^2(G/A_c)'$ where $G(z)/A_c(z)=(b z^p + \lambda z^{p+2})/(mc)$. Since $p+2\in\{1,3\}$ is positive and $p\in\{+1,-1\}$ is nonzero, $\int_1^{v(0)}A_c^2(z^{p+2})'\,\dd z>0$ and $\sgn\int_1^{v(0)}A_c^2(bz^p)'\,\dd z=\sgn(bp)$, giving a unique $\lambda_0\neq 0$ (of sign $-\sgn(bp)$) at which $V(0)=0$. At $(\lambda,\theta)=(\lambda_0,0)$, the integrand of $V$ on $z>1$ has the sign of $bp+(p+2)\lambda_0 z^2$; vanishing of $\int_1^{v(0)}$ forces its unique positive root to lie strictly inside $(1,v(0))$, so the integrand has strict sign $-\sgn(bp)$ throughout $[v(0),v(-1/2)]$, proving $\sgn V(-1/2)=-\sgn(bp)$ at $(\lambda_0,0)$.

\emph{Absolute actions and the two matches.} In scaled variables $z=x/R$, $U=z^2/2$, $t=y/R^m$, running $X_{s,F_R}$ \emph{backwards} in $x>0$ increases $t$ according to $\epsilon\,\dd U/\dd t=\phi_s(U)-t$, where $\Phi_s(z)=\phi_s(z^2/2)=E_\tau(z)+s\tau^g(G(z)+O(\tau^2))$ and $E_\tau(z)=E_c(z)+O(\tau^2)$ collects all even powers of $P(x)-cx^m$. Fix a right wall $z_c>v(-1)$ with $E_c(z_c)<-2$ and a section $U_e=z_e^2/2$ with $z_e=1/2$. Let $H_s^*$ be the scaled backward upper-axis hit from $(U,t)=(0,-1)$, let $h_s(\eta)$ be that from $(U_e,\eta)$ for small $|\eta|$, and put $g_s(\eta)=H_s^*-h_s(\eta)>0$ (positive by arc order since the $(0,-1)$ arc lies near $v(0)^2/2>U_e$ at $t=\eta$). Let $g_s^b>0$ be the analogous gap from $(U,t)=(0,-1/2)$. 
Write $u_s(t)$ for the falling root of $\phi_s(u)=t$ and $M_s$ for
the actual maximum near $z=1$. All constants below are uniform for
$(\lambda,\theta)$ in a fixed compact neighborhood of $(\lambda_0,0)$,
for both sides and for $y$ in any fixed compact interval. The wall
at $z_c$ points inward throughout the passage. On $[0,z_c]$ the
global maximum stays near $1$ by compact convergence and local
nondegeneracy. In particular the compact core cannot create a competing
maximum at these heights.

Fix $0<b_0<E_c(z_e)$. For a launch at $a_0=\tau^m y$ from $U_e$, or
at $a_0=-1/2$ from the axis, the reference from $-1$ is already
$O(\epsilon)$ from the falling root at $a_0$. Indeed, before capture
the velocity towards a fixed root neighborhood is at least $c/\epsilon$
in magnitude. This follows from $\phi_s(U)-t>c$ on the compact
approach outside that neighborhood at the specified launches, including
the initial launch at $-1$. Inside a fixed neighborhood of the regular
falling root, $|U-u_s|$ has upper right derivative bounded by
$-c|U-u_s|/\epsilon+C$. Thus capture in an $O(\epsilon)$ tube takes
$O(\epsilon|\log\epsilon|)$ in $t$. The initial separation from the
reference is bounded above and away from zero. During capture the
secant of $\phi_s$ stays bounded, since either its arguments lie away
from zero or they have a fixed positive separation. On the common tube
it differs from $\phi_s'(u_s(t))$ by $O(\epsilon)$. Integration of the
scalar difference equation up to $b_0$ gives, for
$d_s(a_0)=U_s^*(b_0)-U_s(b_0;a_0)>0$,
\begin{equation}\label{eq:entry-distortion}
 \epsilon\log d_s(a_0)
 =\int_{a_0}^{b_0}\phi_s'(u_s(t))\,\dd t+O(\epsilon|\log\epsilon|).
\end{equation}
The integral is uniformly negative by a fixed amount. The interval
$B_s$ between the extreme section values, including the reference,
therefore has length $O(e^{-c_3/\epsilon})$ for some $c_3>0$.

Every $\xi\in B_s$ tracks the falling branch up to $t=M_s-4\tau^4$.
On the retained branch $-\phi_s'\ge c\tau^2$, $|u_s'|\le C\tau^{-2}$,
and $\phi_s''$ is bounded. A tube of radius $\tau^5$ points inward
because $c\tau^7/\epsilon\gg\tau^{-2}$ for $m\ge6$, and contains all
the arcs at $b_0$. If an arc is still above the axis at $M_s+\tau^4$,
then $\epsilon U'\le-\tau^4$ until it exits, in an additional
$O(\epsilon/\tau^4)$ of height. Its exit $K_s(\xi)$ consequently lies
within $O(\tau^4)$ of $M_s$.

Choose a fixed $0<u_c<U_e$ so small that on the final strip
$0\le U\le 2u_c$ near $M_s$ one has $t-\phi_s(U)\ge\nu>0$.
Up to $u_c$, $f=\phi_s(U)-t$ has uniformly bounded first two
$U$ derivatives. Inside the retained tube $f_U<0$, also between
its arcs. The remaining interval before crossing $u_c$ has length
$O(\tau^4)$, so the quantity $D$ of Lemma~\ref{lem:terminal-distortion}
is $O(\tau^4)$. The inverse strip has $f_t=-1$, $f_{tt}=0$.
Thus the lemma gives
\begin{equation}\label{eq:cap-osc}
 \operatorname{osc}_{B_s}(\epsilon\log K_s')=O(e^{-c_4/\epsilon})
 \qquad(0<c_4<c_3)
\end{equation}
for small $\tau$. For the absolute size we also have
\begin{equation}\label{eq:Ks-deriv}
 K_s'(\xi)=\frac{\epsilon}{K_s(\xi)}
 \exp\!\left(\epsilon^{-1}\int_{b_0}^{K_s(\xi)}
                           \phi_s'(U(t;\xi))\,\dd t\right).
\end{equation}
Indeed differentiate first at a section $U=u>0$ and then let
$u\downarrow0$. On the last strip $\dd t/\dd U=\epsilon/(\phi_s(U)-t)$,
and $\int_0^{u_c}|\phi_s'|\,\dd U$ is uniformly finite, including
the possible $U^{-1/2}$ term. Hence this part of the integral costs
$O(\epsilon)$, the section factor tends to $\epsilon/K_s$, and the
inverse-chart derivative converges by the proof of the lemma.
On the retained tube the action error is $O(\tau^5)$, and both the
omitted root integral and the actual fold passage above $u_c$ cost
$O(\tau^4)$. Together with~\eqref{eq:entry-distortion} and a secant
of $K_s$, this gives the absolute falling action to $O(\tau^4+
\epsilon|\log\epsilon|)$.
Let $\widehat z_e$ and $\widehat z_s$ be the respective cap maximizers
of $E_\tau$ and $\Phi_s$, and $v_e(Y),v_s(Y)$ their falling roots at
height $Y$. Set $M_\tau=E_\tau(\widehat z_e)$,
$A_\tau=E_\tau'/z$ and
$J_{o,\tau}(Y)=\int_{\widehat z_e}^{v_e(Y)}E_\tau'(z)^2\,\dd z/z$.
For $Y$ near $0$ or $-1/2$, the upper endpoint of the positive action
moves by $-s\tau^gG(v_e)/E_\tau'(v_e)+O(\tau^{2g}+\tau^{g+2})$.
Its lower endpoint contributes zero at first order because
$E_\tau'(\widehat z_e)=0$. The coefficient of $s\tau^g$ is consequently
\[
 2\int_{\widehat z_e}^{v_e(Y)} A_\tau G'\,\dd z-A_\tau(v_e(Y))G(v_e(Y))
 =\int_{\widehat z_e}^{v_e(Y)}(A_\tau G'-A_\tau'G)\,\dd z
 = V(Y)+O(\tau^2).
\]
De Maesschalck and Huzak compute this first variation near a center
\cite[Proposition~1]{DMH2015}. For $F=F_e+\delta F_o$ with $F_e$ even,
$F_o$ odd and $f_e=F_e'/z$, their integral from the right endpoint to
the left endpoint at the same height is $I=2\delta I_1+O(\delta^2)$, where
$I_1(x)=\int_0^x(f_e'F_o-f_eF_o')\,\dd z$.
Our positive action runs from the maximum to the falling root. With
$(f_e,F_o)=(A_\tau,G)$ its coefficient is the negative increment
$I_1(\widehat z_e)-I_1(v_e(Y))$. The action enters $\epsilon\log g_s$
with a minus sign. We keep the actual even profile on this interval to control the
remainder uniformly in the matching parameters.
The maximum height changes by $s\tau^g\theta+O(\tau^{2g}+\tau^{g+2})$
since $\widehat z_e=1+O(\tau^2)$ and $G(1)=\theta$. Replacing the
bounded physical launch level $\tau^m y$ by zero in the absolute
action costs $O(\tau^m)$. For $g\in\{1,3\}$ all these remainders
and the $O(\tau^4)$ passage error are $O(\tau^{g+1})$. We obtain
\begin{align}
  H_s^* &= M_\tau + s\tau^g\theta + O(\tau^{g+1}),\nonumber\\
  \epsilon\log g_s(\tau^m y) &= -J_{o,\tau}(0) - s\tau^g V(0) + O(\tau^{g+1}),\label{eq:cap-coarse}\\
  \epsilon\log g_s^b &= -J_{o,\tau}(-1/2) - s\tau^g V(-1/2) + O(\tau^{g+1}).\nonumber
\end{align}

On the rectangle $\mathcal{Q}_\tau=[\lambda_0-\rho_0,\lambda_0+\rho_0]\times[-\sqrt\tau,\sqrt\tau]$ ($0<\rho_0<|\lambda_0|/2$ fixed), $H_+^*-H_-^*=2\tau^g\theta+O(\tau^{g+1})$ has opposite signs on the $\theta=\pm\sqrt\tau$ faces, and $\epsilon\log(g_+(\tau^m\Delta)/g_-(0))=-2\tau^g V(0)+O(\tau^{g+1})$ has opposite signs on the $\lambda=\lambda_0\pm\rho_0$ faces because $\partial_\lambda V(0)>0$. Orienting the continuous residual vector $F_{\mathrm{res}}(\lambda,\theta)$ outward on the boundary of $\mathcal{Q}_\tau$ and applying Brouwer's fixed-point theorem to $\operatorname{proj}_{\mathcal{Q}_\tau}\big((\lambda,\theta)-\nu F_{\mathrm{res}}(\lambda,\theta)\big)$ ($\nu>0$ small) yields an interior root $(\lambda_R,\theta_R)\in\operatorname{int}\mathcal{Q}_\tau$ solving simultaneously
\begin{equation}\label{eq:cap-matches}
  H_+^* = H_-^* =: H_*,\qquad g_+(\tau^m\Delta) = g_-(0),
\end{equation}
with $\theta_R=O(\tau)$ and $\lambda_R=\lambda_0+O(\tau)\neq 0$. Thus $F_R$ has exact degree $\max(m+2,j+4)$.

\emph{Matching a finite list at physical resolution.}
The error in~\eqref{eq:cap-coarse} is of order $\tau^{g+1}$, whereas a
bounded difference of physical ordinates has scaled size $\tau^m$.
To resolve that smaller scale, split propagation at $t=b_0$ and use
the same terminal map $K_s$ for every launch on a given side.
Equations~\eqref{eq:entry-distortion} and~\eqref{eq:cap-osc} apply on the
entire interval between the launches and the reference. By
Lemma~\ref{lem:terminal-distortion}, comparing the two positive terminal
gaps costs at most $O(e^{-c_4/\epsilon})$ in their scaled log ratio. Hence
\begin{align}
 \frac{\epsilon}{\tau^m}\log\frac{g_s(\tau^m y)}{g_s(\tau^m y_0)}
 &= -\tau^{-m}\int_{\tau^m y_0}^{\tau^m y}\phi_s'(u_s(t))\,\dd t
       +O\!\left(\frac{\epsilon|\log\epsilon|+e^{-c_4/\epsilon}}{\tau^m}\right)
       \nonumber\\
 &= -A_c(v(0))(y-y_0)+o(1).\label{eq:cap-fine}
\end{align}
The error is uniform for $y,y_0$ in compact sets,
since $\epsilon|\log\epsilon|/\tau^m=(2m-2)\tau^{m-2}|\log\tau|\to0$.
At $v(0)^2=(m+2)/m$ the coefficient is
$-A_c(v(0))=2c((m+2)/m)^{(m-2)/2}>0$.
At $U_e$ the backward field points right, so earlier launches lie
farther right and $h_s(\tau^m y)$ decreases strictly with $y$.
Using~\eqref{eq:cap-matches} and~\eqref{eq:cap-fine} with base ordinates
$\Delta$ on the positive side and zero on the negative side brackets
the opposite ordinate between $y-\Delta-\rho$ and $y-\Delta+\rho$ for
every fixed $\rho>0$ and all sufficiently small $\tau$.
Strict monotonicity gives a unique $y_-=y-\Delta+o(1)$ with
$h_-(\tau^m y_-)=h_+(\tau^m y)$.

To continue these matched arcs inward, take $Z_R=R z_e=R/2$.
On $X_0\le x\le Z_R$, division of~\eqref{eq:cap-poly} by $x^m$ gives
$F_R(sx)/x^m=c(1-\tfrac{m}{m+2}(x/R)^2)+O(x^{-1})$ uniformly
for the bounded matching parameters. Thus $F_R(sx)\ge c_0x^m$ there
for fixed $c_0>0$ and sufficiently large fixed $X_0$.
We already have $F_R\to P$ on compact sets in $C^2$.
Corollary~\ref{cor:transverse-entry} applies at $y_i$ on the positive
side and at $y_i-\Delta$ on the negative side. Its local uniformity
and continuity allow the negative launch to move to $y_-$.
At their common physical upper height $R^m h_+(\tau^m y_i)$ the
half-return difference therefore converges to
$I_{+,P}(y_i)-I_{-,P}(y_i-\Delta)$. The list is finite and the signs
are strict, so all $r$ changes persist together, in increasing upper height.

\emph{The last two heights.} Since $J_o(0)<J_o(-1/2)$, both backward upper hits $H_*-g_\pm^b$ from $(0,-R^m/2)$ lie strictly above all section upper hits $h_+(\tau^m y_i)$ and strictly below $H_*$. By~\eqref{eq:cap-coarse}, $\epsilon\log(g_+^b/g_-^b)=-2\tau^g V(-1/2)+O(\tau^{g+1})$ has strict sign $-\sgn V(-1/2)=\sgn(bp)$. If $g_+^b>g_-^b$, the plus upper hit $H_*-g_+^b$ is smaller than the minus upper hit $h_{\mathrm{dir}}:=H_*-g_-^b$, so monotonicity of $w_+$ on the origin-free strip $[R^m h_+(\tau^m y_0),R^m H_*]$ (Lemma~\ref{lem:half-returns}(iii), as both endpoints return transversely to $\Sigma_-$) gives $w_+(R^m h_{\mathrm{dir}})>R^m/2=w_-(R^m h_{\mathrm{dir}})$, i.e., $\D(R^m h_{\mathrm{dir}})>0$; if $g_+^b<g_-^b$, reversing roles gives $\D<0$ at $R^m(H_*-g_+^b)$. Thus a sample of strict sign $\sgn(bp)$ precedes $R^m H_*$, where $\D(R^m H_*)=0$. The entrance list has $r$ changes and ends with sign $\sigma$.
The direct sample adds one change exactly when $\sigma\ne\sgn(bp)$.
Keep it as the representative of the last run if the signs agree.
Detuning the matched reference to $-\sgn(bp)$ then adds one more
change in either case. Proposition~\ref{prop:sard}(i) on
$[R^m h_+(\tau^m y_0),R^m H_*]$ realizes all these brackets at one
regular value. The $q$ hyperbolic cycles in $K_0$ persist, disjoint
from the new heights $\asymp R^m E_c(1)$, giving~\eqref{eq:cap-count}.
\end{proof}

\begin{example}[Repelling passage with an oblique exit]\label{ex:repelling}
To see the role of the margin $A-D$ in Lemma~\ref{lem:terminal-distortion},
consider a time-dependent scalar equation, or equivalently its planar
extension $\dot t=1$. Write $U=2+x$. Until $t=0$ take
$\epsilon x'=-x$ and launch $x=q\in(0,1)$ at $t=-A+\delta y$,
where $A>0$, $\delta\downarrow0$ and $\epsilon=o(\delta)$.
The reference is $x=0$. On $0\le t\le1$ use
$\epsilon x'=D x(1-x)$ with $D>0$. After $t=1$ take
$\epsilon U'=-1$ and stop on the oblique line $U=\alpha+\gamma(t-1)$,
with $0<\alpha<1$ and fixed $\gamma$. At $t=0$ the entering
distance is $d_y=q e^{(-A+\delta y)/\epsilon}$. The exact logistic
flow and the final straight flight give the terminal time gap
\begin{equation}\label{eq:logistic-gap}
 g_y=\frac{\epsilon}{1+\epsilon\gamma}\,
        \frac{e^{D/\epsilon}d_y}{1+(e^{D/\epsilon}-1)d_y}.
\end{equation}
On the last flight $U-\alpha-\gamma(t-1)$ straightens the section and
has uniformly negative fast velocity. On the repelling interval
$(f_U)_+\le D$. Thus the lemma applies with a strictly positive
margin for bounded $y$ when $D<A$.
More precisely, for fixed $A,D$ and $y$,
\[
 \lim\frac{\epsilon}{\delta}\log(g_y/g_0)=
 \begin{cases}
 y,&D<A,\\
 \min(y,0),&D=A,\\
 0,&D>A .
 \end{cases}
\]
Indeed the denominator in~\eqref{eq:logistic-gap} is
$1+q(1-e^{-D/\epsilon})e^{(D-A+\delta y)/\epsilon}$.
At equality its logarithm, multiplied by $\epsilon/\delta$, tends to
$\max(y,0)$. Above equality the denominators cancel the entire input
slope. At $D=A$ a two-sided window of launches includes inputs wider
than $q e^{-A/\epsilon}$ by the factor $e^{\delta y/\epsilon}$ for $y>0$.
On the smaller bundle $0\le d\le d_0=q e^{-A/\epsilon}$, differentiating
the passage in~\eqref{eq:logistic-gap} at $D=A$ gives
\[
  K'(0)/K'(d_0)=\big(1+q(1-e^{-A/\epsilon})\big)^2
       \longrightarrow(1+q)^2.
\]
Thus the oscillation of
$\epsilon\log K'$ at $D=A$ is
$2\epsilon\log(1+q(1-e^{-A/\epsilon}))$. This is of order $\epsilon$,
hence $o(\delta)$ on the smaller bundle. A nonnegative margin alone
cannot guarantee exponential distortion control.

For each sufficiently small $\epsilon>0$ a smooth planar field realizes
the same gaps. Replace the two rates by
smooth weights with disjoint time supports, equal to $1$ near the
launch for the contracting rate, with integrals $A$ on $[-A,0]$
and $D$ on $[0,1]$, respectively. After $t=1$ use fast velocity
$-\chi((t-1)/\epsilon)$, with $\chi$ smoothly increasing from zero
to one on $[0,1]$, constant outside. The common displacement during
this ramp is at most one. Since $\alpha<1$, the ramp stays a fixed
distance from the exit for small $\epsilon$. All hits occur on the
following constant-speed flight and still satisfy~\eqref{eq:logistic-gap}.
\end{example}

